\documentclass[pdflatex,sn-mathphys-num]{sn-jnl}
% Official Springer Nature template, December 2024 package; class unmodified.
\usepackage{amsmath,amssymb,amsthm,mathtools,booktabs,array}
\usepackage{xcolor}
\usepackage[T1]{fontenc}
\hypersetup{hidelinks,bookmarksdepth=2}
\theoremstyle{thmstyleone}
\newtheorem{theorem}{Theorem}
\newtheorem{lemma}[theorem]{Lemma}
\newtheorem{corollary}[theorem]{Corollary}
\theoremstyle{thmstyletwo}\newtheorem{remark}[theorem]{Remark}
\newcommand{\R}{\mathbb R}
\newcommand{\K}{\mathcal K}
\newcommand{\PSD}{\mathbb S_+}
\newcommand{\sxdeg}{\operatorname{sxdeg}}
\newcommand{\esssup}{\operatorname*{ess\,sup}}
\newcommand{\diag}{\operatorname{diag}}
\newcommand{\spn}{\operatorname{span}}
\newcommand{\ip}[2]{\langle #1,#2\rangle}
\allowdisplaybreaks
\setlength{\emergencystretch}{2em}
\raggedbottom
\AddToHook{env/thebibliography/after}{\clearpage}
\begin{document}
\title[Minimum-Peak Control via Exact Semidefinite Lifts]{Minimum-Peak Control
of Fifth-Order Integrators via Exact Semidefinite Lifts}
\author*[1]{\fnm{Zihan} \sur{Yang}}\email{768216265@qq.com}
\affil[1]{\orgdiv{School of Information and Communication Engineering},
\orgname{University of Electronic Science and Technology of China},
\orgaddress{\city{Chengdu}, \country{China}}}

\abstract{
We study minimum-peak control of fifth-order integrator chains with known
initial states and linear state constraints at prescribed times. An exact
finite model allows these constraints to be assembled while leaving the
input waveform unrestricted within each time segment. We establish such a
semidefinite programming (SDP) formulation without an outer input-time mesh.
Its main ingredient is an explicit affine
lift of the complete fifth-order bounded-density moment body, using 34
auxiliary scalars and 21 positive semidefinite (PSD) blocks of order at most four.
A domain-weighted first-order identity and finite witnesses cover singular
boundaries; homogenization and segment concatenation preserve equivalence,
including zero peak. An application of existing cone-complexity theory also
excludes finite affine second-order-cone lifts of the full reach set in
orders five and above. The fifth-order semidefinite extension degree lies
between three and four. Synthesis tests compare the SDP implementation with
a continuous-dual solver using shared rational recovery and certification.
Both return valid bounds on all tested tasks, although some bounds miss
the requested accuracy. The SDP is slower on every jointly successful task,
and feedback replay activates no SDP plan under a half-second deadline.
The results support an exact finite interface for auditable offline synthesis;
certification is shared by both methods, and no real-time advantage is claimed.}

\keywords{Minimum-peak control, integrator chains, semidefinite lifts,
bounded-density moments, extension degree}
\pacs[Mathematics Subject Classification (2020)]{93B03, 90C22, 44A60}
\maketitle

\section{Introduction}
Minimum-peak control with fixed-time linear state constraints is an
infinite-dimensional convex problem over measurable inputs. This paper
constructs an equivalent finite affine semidefinite programming (SDP) model
for fourth- and fifth-order integrator chains with componentwise input bounds.
The contribution concerns exact representation, including singular boundaries
and zero peak, rather than a new control objective or a faster planner.

Integrator reach-set geometry and support-function methods are well studied
\cite{hh,intersection}. Boundary parameterizations do not by themselves provide
affine conic constraints in the original state coordinates. The SOCP intersection
tests in \cite{intersection} address a support-function decision problem with
numerical quadrature, not an exact finite lift of the complete reach set.
Likewise, classical
Markov-moment transforms \cite{gr} are nonlinear coordinate descriptions, and
full-sequence domination conditions (Lasserre, 2007 preprint,
\href{https://arxiv.org/abs/math/0607463v2}{arXiv:math/0607463v2})
quantify over all polynomial degrees. Lasserre's later density-existence
hierarchy \cite{lasserredensity} likewise does not by itself identify a fixed
finite exact level for the five-moment body. First-order polynomial and rational lifting
\cite{nie,nierational} and closure of spectrahedral shadows \cite{gn} supply
established tools. The specific construction below supplies the weighted
certificate, common coefficient space and finite boundary witnesses needed
for an explicit fifth-order realization.

Minimum-maximum-derivative interpolation and perfect splines also have
classical precedents \cite{mcclure}. Our fixed-time state constraints differ
from continuous collision-free polynomial planning with integral-squared costs
and nonconvex relaxations (Graesdal et al., 2026 preprint,
\href{https://arxiv.org/abs/2606.14063v1}{arXiv:2606.14063v1}). No vehicle dynamics, unknown arrival
times or inter-knot obstacle constraints are included.

The main results are:
\begin{enumerate}
\item An affine lift of the complete five-moment body with 34 auxiliary
scalars and 21 positive-semidefinite (PSD) blocks of maximum order four,
including singular boundaries (Theorem~\ref{thm:five-lift}).
\item Exact homogenization and segment concatenation, including zero peak,
yield a finite SDP for fixed-time minimum-peak synthesis
(Corollary~\ref{cor:five-synthesis}).
\item A normal-cone argument using established lift duality \cite{gpt} and
cone-complexity theory \cite{averkov} gives
$3\le\sxdeg(\K_5)\le4$ and excludes finite second-order-cone (SOC) lifts
of full reach sets in orders at least five (Theorem~\ref{thm:lower}).
\end{enumerate}
The fifth-order case is a low-order setting where an explicit finite SDP
and a non-SOC lower bound coexist. The exact block order three versus four
and SOC representability of the fourth-order body remain open here.
The lower bound concerns representations, not the speed of specialized
algorithms. The continuous-dual reference accepts the same linear task rows
and shares the rational recovery and checking procedure; neither capability
is exclusive to SDP.

Section~\ref{sec:control-problem} defines the system and moment coordinates.
Section~\ref{sec:theory} gives the lifts, representation bound and synthesis
equivalence. Section~\ref{sec:validation} tests their numerical realization
and implementation limits, retaining accuracy failures and missed deadlines.

\section{System modeling}\label{sec:control-problem}
\subsection{Dynamics, constraints and objective}
Consider $A$ independent chains of length $n\in\{4,5\}$, with state
$x_{a,j}$ for axis $a=1,\ldots,A$ and derivative index $j=0,\ldots,n-1$:
\[
 \dot x_{a,j}=x_{a,j+1}\quad (0\le j<n-1),\qquad
 \dot x_{a,n-1}=u_a.
\]
Stack all states into $x$. Initial states and the horizon $T>0$ are known.
Given finitely many fixed times $t_i$, row vectors $w_i$, targets $b_i$ and
tolerances $\epsilon_i$, require
\begin{equation}\label{eq:waypoints}
 |w_i^\top x(t_i)-b_i|\le\epsilon_i,\qquad 0<t_i\le T,
 \quad\epsilon_i\ge0.
\end{equation}
Rows may couple axes and derivatives. Minimize
$J=\esssup_{0\le t\le T}\max_a|u_a(t)|$ over
$u\in L^\infty([0,T];\R^A)$.
The bound is componentwise, and no path constraint between the listed
times is imposed.

The finite model must preserve all feasible peak values, including zero.
Partitioning at row times leaves the input waveform unrestricted inside each
segment; Section~\ref{sec:synthesis} proves equivalence using complete moment
bodies rather than an input-time grid.

The scope is offline synthesis for independent chains with fixed observation
times. Unknown arrival times, coupled input balls, continuous obstacle
avoidance and nonlinear plant dynamics are excluded. Relating the highest
derivative peak to physical actuator loads requires a separate plant model.

\subsection{Moment coordinates and affine lifts}\label{sec:moments}
For an integer $d\ge1$, let
\begin{equation}\label{eq:body}
 \K_d=\left\{y\in\R^d:y_j=\int_0^1 t^j\rho(t)\,dt,
 \quad 0\le\rho\le1\ \text{a.e.},\quad j=0,\ldots,d-1\right\}.
\end{equation}
The order interval of densities is weak-star compact in $L^\infty$, so
$\K_d$ is compact and convex. With $v_j=1/(j+1)$ and $c=v/2$, its centered
support function is
\begin{equation}\label{eq:support}
 h_{\K_d-c}(\lambda)=\frac12\int_0^1
 \left|\sum_{j=0}^{d-1}\lambda_jt^j\right|\,dt.
\end{equation}
It is positive for every nonzero $\lambda$, hence $\K_d$ is full-dimensional.

We first fix a candidate peak $J>0$ to describe its reach set; the
minimum-peak formulation subsequently optimizes this same parameter.
Consider $\dot x_1=x_2,\ldots,\dot x_{d-1}=x_d,\dot x_d=u$,
with known $x(0)$, $T>0$ and $|u|\le J$, where $J>0$ is fixed.
Set $u(Tt)=J(2\rho(t)-1)$. The input contribution to $x_{d-k}(T)$ is
\begin{equation}\label{eq:affine}
 \frac{JT^{k+1}}{k!}\left[
 2\sum_{j=0}^k(-1)^j\binom{k}{j}y_j-\frac1{k+1}\right],
 \qquad k=0,\ldots,d-1.
\end{equation}
This is an invertible triangular affine map: its diagonal is
$2JT^{k+1}(-1)^k/k!$. Adding the known zero-input flow therefore gives an
affine equivalence between the reach set and $\K_d$. For $A$ independent
chains with componentwise input bounds, the full reach set is the Cartesian
product of their $d$-dimensional single-axis reach sets. Coupled Euclidean
input balls do not follow from this product argument.

A convex set has an affine SDP lift if it is the projection of a set defined
by finitely many linear matrix inequalities and linear equations. Following
\cite{averkov}, $\sxdeg(C)$ is the least possible maximum PSD block order,
allowing arbitrarily many finite blocks and auxiliaries; it is infinity
if no finite lift exists. A nonlinear change of the output coordinates is
not an affine lift.

\section{Theoretical derivation}\label{sec:theory}
We first construct the four-moment prefix, then the complete fifth-order
lift and its coefficient interface. A lower bound distinguishes its cone
class. Homogenization and concatenation finally connect the representation
to the system model.

\subsection{Four-moment prefix lift}\label{sec:four}
Write $y=(m,z,q,r)$ and
$\bar y=(1-m,\tfrac12-z,\tfrac13-q,\tfrac14-r)$. Define
\begin{align}
 D&=z-m^2/2,& E&=m-z-m^2/2,\notag\\
 B&=q-mz/2-m^3/12,& F_0&=z^2/2+m^2z/4,\notag\\
 G(y)&=\begin{pmatrix}D&B\\B&r-F_0\end{pmatrix}.\label{eq:nonlinear}
\end{align}

\begin{lemma}[Classical boundary condition in convenient coordinates]\label{lem:G}
 $y\in\K_4$ if and only if $G(y)\succeq0$ and $G(\bar y)\succeq0$.
\end{lemma}
\begin{proof}
Section~\ref{app:boundary} gives an interval-extremizer derivation,
including all zero-denominator cases. To see why the prefix cannot be
silently omitted, when $0<m<1$ and $D,E>0$, put
\[
 q_- = z^2/m+m^3/12,\qquad
 q_+ =1/3-(1/2-z)^2/(1-m)-(1-m)^3/12.
\]
The lower fourth-moment boundary is
$L(y)=F_0+B^2/D$. Direct expansion gives
\begin{align}
 q_+-q_-&=\frac{DE}{m(1-m)},\notag\\
 L(y)+L(\bar y)-1/4
 &=\frac{m(1-m)(q-q_-)(q-q_+)}{DE}.\label{eq:prefix}
\end{align}
The two PSD conditions imply $D,E\ge0$ and $0\le m\le1$.
Their Schur complements and \eqref{eq:prefix} imply $q_-\le q\le q_+$,
the feasible three-moment prefix. The lower and complementary upper
fourth-moment bounds are then necessary and sufficient.
If $0<m<1$ and $D=0$, PSD forces $B=0$, hence
$z=m^2/2,q=m^3/3$; the two last diagonal/Schur bounds force $r=m^4/4$.
The case $E=0$ follows by complementation. If $m=0$, the matrices force
$z=q=r=0$; $m=1$ follows by complementation. Conversely each feasible
density satisfies the same extremal bounds.
\end{proof}

The matrices in Lemma~\ref{lem:G} are nonlinear in $y$. To obtain an
affine lift, introduce six auxiliary scalars $a,b,c,d,e,f$ and the pencils
\begin{align}
 H&=\begin{pmatrix}1&m&a\\m&a&b\\a&b&c\end{pmatrix},\label{eq:H}\\
 J_\sigma&=\begin{pmatrix}
 1&m&z+\sigma a/6\\m&a&d+\sigma b/6\\
 z+\sigma a/6&d+\sigma b/6&f+\sigma e/3+c/36
 \end{pmatrix},\qquad \sigma\in\{-1,1\},\label{eq:J}\\
 T_-&=\begin{pmatrix}z-a/2&d-b/2\\d-b/2&e-c/2\end{pmatrix},\notag\\
 T_+&=\begin{pmatrix}m-z-a/2&a-d-b/2\\a-d-b/2&b-e-c/2\end{pmatrix},\label{eq:T}\\
 Q_-&=\begin{pmatrix}z-a/2&q-d/2-b/12\\q-d/2-b/12&r-f/2-e/4\end{pmatrix},\notag\\
 Q_+&=\begin{pmatrix}m-z-a/2&B_+\\ B_+&C_+\end{pmatrix},\label{eq:Q}\\
 B_+&=-a/4+b/12-d/2+m/2-q+z/2,\notag\\
 C_+&=-a/8-d/2+e/4-f/2+m/4-r+3z/4.\notag
\end{align}

\begin{theorem}[Four-moment affine lift]\label{thm:lift}
 A vector $y\in\R^4$ belongs to $\K_4$ if and only if there exist
 $a,b,c,d,e,f\in\R$ making all seven matrices
 $H,J_+,J_-,T_-,T_+,Q_-,Q_+$ positive semidefinite.
 In particular, $\sxdeg(\K_4)\le3$.
\end{theorem}

\begin{proof}
For $g\in\R$, define
\[
 F_g(s,t)=t^2/2+s^2t/4-g(s^3/12+st/2)
             +g^2(s^2/8-t/4).
\]
Set $u=s-M$, $w=t-Z$, $D(s,t)=t-s^2/2$, $D_0=Z-M^2/2$.
The following polynomial identity follows by expansion, or by integrating
the Hessian along the line segment from $(M,Z)$ to $(s,t)$:
\begin{align}
 &F_g(s,t)-F_g(M,Z)-\nabla F_g(M,Z)^\top(u,w)\notag\\
 &\quad=\frac12\left[w+\frac{M-g}{2}u+\frac{u^2}{6}\right]^2
       +\frac{D(s,t)u^2}{12}+\frac{D_0u^2}{6}+\frac{u^4}{36}.
       \label{eq:remainder}
\end{align}
This is a domain-weighted first-order certificate, not an invocation of
global SOS-convexity.

Given a feasible lifted tuple, define a linear functional $\Lambda$ by the
assignments
\[
\begin{array}{c|rrrrrrrrr}
g&1&s&t&s^2&s^3&s^4&st&s^2t&t^2\\
\Lambda(g)&1&m&z&a&b&c&d&e&f
\end{array}
\]
No representing measure for $\Lambda$ is assumed. The matrices $H$ and $J_+$
certify nonnegativity of $\Lambda$
on squares in $\spn\{1,s,s^2\}$ and $\spn\{1,s,t+s^2/6\}$, respectively.
The matrix $T_-$ certifies $\Lambda[D(s,t)h^2]\ge0$ for
$h\in\spn\{1,s\}$.

From $H\succeq0$ we obtain $a\ge m^2$. The leading entries of $T_\pm$
give $a\le m$ and $z,m-z\ge a/2$. Thus $0\le m\le1$ and
$D_0=z-m^2/2\ge0$. Apply $\Lambda$ to \eqref{eq:remainder} with
$(M,Z)=(m,z)$. The linear Taylor term vanishes, and each term on the right
is nonnegative by the three named matrices and $D_0\ge0$. Hence
\begin{equation}\label{eq:jensen}
 \Lambda(F_g)\ge F_g(m,z),\qquad g\in\R.
\end{equation}
The quadratic form of $Q_-$ at $(-g/2,1)$ yields
\begin{equation}\label{eq:allg}
 r\ge gq+\Lambda(F_g)\ge gq+F_g(m,z),\qquad g\in\R.
\end{equation}
If $D_0>0$, maximizing the last expression over $g$ gives
$r\ge F_0+B^2/D_0$. If $D_0=0$, its finiteness forces $B=0$ and
$r\ge F_0$. In both cases $G(y)\succeq0$.

For the complementary coordinates $\bar s=1-s$, $\bar t=1/2-t$, one has
\begin{align*}
 \spn\{1,\bar s,\bar s^2\}&=\spn\{1,s,s^2\},\\
 \spn\{1,\bar s,\bar t+\bar s^2/6\}&=\spn\{1,s,t-s^2/6\},\\
 D(\bar s,\bar t)&=s-t-s^2/2.
\end{align*}
Thus $H,J_-,T_+$ give the same argument. In detail, their transformed
auxiliaries are
\begin{align*}
 \bar a&=1-2m+a,&\bar b&=1-3m+3a-b,\\
 \bar c&=1-4m+6a-4b+c,&\bar d&=1/2-z-m/2+d,\\
 \bar e&=1/2-m+a/2-z+2d-e,&\bar f&=1/4-z+f.
\end{align*}
Substitution in $Q_-$ gives exactly $Q_+$. Consequently
$G(\bar y)\succeq0$. Lemma~\ref{lem:G} proves sufficiency, including the
singular faces.

For necessity take any $y\in\K_4$ and choose
\[
 (a,b,c,d,e,f)=(m^2,m^3,m^4,mz,m^2z,z^2).
\]
Then $H$ and $J_\sigma$ are outer products,
$T_-=D(1,m)^\top(1,m)$ and $T_+=E(1,m)^\top(1,m)$.
The last two matrices are $G(y)$ and $G(\bar y)$, so all constraints hold.
\end{proof}

\begin{remark}[Attribution and the role of the auxiliaries]
The boundary condition is classical; its affine lift is the constructive object.
For example, the first four coefficients in
$1-\exp[-(m w+z w^2+q w^3+r w^4+\cdots)]$ are
\[
 m,\quad z-m^2/2,\quad q-mz+m^3/6,\quad
 r-mq-z^2/2+m^2z/2-m^4/24.
\]
If the last three are denoted $a_1,a_2,a_3$, then
\[
 \begin{pmatrix}1&0\\m/2&1\end{pmatrix}
 \begin{pmatrix}a_1&a_2\\a_2&a_3\end{pmatrix}
 \begin{pmatrix}1&m/2\\0&1\end{pmatrix}=G(y).
\]
This congruence links Lemma~\ref{lem:G} to the classical transform.
The auxiliaries are coordinate-polynomial functional values, not extra time
moments. Standard first-order lifting \cite{nie} motivates the construction.
A generic polynomial-coordinate lift, defined below, has block orders
$6,3,3,2,2$ and twelve auxiliaries with the same projection;
Section~\ref{sec:matched} compares it with the seven-block model.
Section~\ref{app:rational} records a separate
rational-lifting existence route, which does not supply the seven-block bound.
\end{remark}

\noindent\textit{Full definition of the generic baseline.}
Let $\Lambda(s^it^j)=\mu_{ij}$ for $i+j\le4$, with
$\mu_{00}=1$, $\mu_{10}=m$ and $\mu_{01}=z$. The other twelve values are
free auxiliary scalars; $q,r$ remain output coordinates. Set
\[
 V_2=(1,s,t,s^2,st,t^2)^\top,\qquad V_1=(1,s,t)^\top,
 \qquad D=t-s^2/2,\quad E=s-t-s^2/2.
\]
The five affine PSD constraints are
\[
 \Lambda(V_2V_2^\top)\succeq0,\quad
 \Lambda(DV_1V_1^\top)\succeq0,\quad
 \Lambda(EV_1V_1^\top)\succeq0,\quad Q_-\succeq0,\quad Q_+\succeq0,
\]
where $Q_\pm$ are exactly \eqref{eq:Q}, with
$(a,b,c,d,e,f)=(\mu_{20},\mu_{30},\mu_{40},\mu_{11},\mu_{21},\mu_{02})$.
All functional evaluations are entrywise on polynomials of degree at most four.
Restricting the first Gram block to $(1,s,s^2)$ and
$(1,s,t\pm s^2/6)$ gives $H,J_\pm$; restricting the two localizing
blocks to $(1,s)$ gives $T_\pm$. Hence generic feasibility implies
the seven-block constraints. Conversely, for every $y\in\K_4$,
$\mu_{ij}=m^iz^j$ gives weighted outer products and the two physical
$G$ matrices, including boundary points. This proves equality of the two
output projections. Homogenize the already linearized pencils by replacing
the constant one by $J$ and using homogeneous output and auxiliary values.
The same restrictions hold at $J=0$, so Lemma~\ref{lem:homogeneous}
excludes additional output directions there.

\subsubsection{Interval-extremizer construction}\label{app:boundary}
Here we provide the classical interval-extremizer argument in the coordinates
used above. For fixed $0<m<1$ and $m^2/2\le z\le m-m^2/2$, the density
$\mathbf1_{[a,a+m]}$ with $a=z/m-m/2$ minimizes the second moment, giving
$q_-=z^2/m+m^3/12$. Indeed $(t-a)(t-a-m)$ is nonpositive on that interval
and nonnegative outside; multiplying by the difference from this density
and integrating gives the lower bound because the first two moments cancel.
Apply the same argument to the complementary density for $q_+$.
Convex combinations of the two extremizers realize every intermediate $q$.

For a feasible prefix $(m,z,q)$ put $D=z-m^2/2$ and $Q=q-m^3/3$.
The prefix bounds become
\[
 0\le D\le m(1-m),\qquad
 mD+D^2/m\le Q\le(m+1)D-D^2/(1-m).
\]
If $D>0$, set $b=Q/D-m$ and
\[
 \Delta=(m-b)^2+4D,\qquad
 a=(m+b-\sqrt\Delta)/2,\qquad
 c=(m+b+\sqrt\Delta)/2.
\]
The bounds imply $D/m\le b\le1-D/(1-m)$, whence
$0\le a\le b\le c\le1$. In particular $mb\ge D$ and
$(1-m)(1-b)\ge D$ imply the endpoint inequalities by comparing
$\Delta$ with $(m+b)^2$ and $(2-m-b)^2$.
The density $\rho_- =\mathbf1_{[0,a]}+\mathbf1_{[b,c]}$ realizes the prefix:
\begin{align*}
 a+c-b&=m,\\
 (a^2+c^2-b^2)/2&=m^2/2+D=z,\\
 (a^3+c^3-b^3)/3&=m^3/3+D(m+b)=q.
\end{align*}
For any competing density with this prefix, the product
$W(t)=(t-a)(t-b)(t-c)$ satisfies $W(\rho-\rho_-)\ge0$ pointwise.
Its lower-degree terms integrate to zero. Therefore the fourth moment is
at least that of $\rho_-$, namely
\[
 L=Q^2/D-mQ+m^2D+D^2/2+m^4/4=F_0+B^2/D.
\]
The complementary construction attains $1/4-L(\bar y)$ as the upper bound.
Their convex combinations fill the interval of feasible fourth moments.
If $D=0$, the sign argument with $t-m$ forces
$\rho=\mathbf1_{[0,m]}$ almost everywhere, hence
$(q,r)=(m^3/3,m^4/4)$. The cases $m=0$ and $m=1$ are the zero and full
densities. This proves the necessary boundary formulas without assuming
nondegeneracy of a representing density.

\subsubsection{Relationship to rational first-order lifting}\label{app:rational}
This alternative fourth-order route clarifies the relationship to existing tools.
Use $x=(s,t,v)$ and a reference $X=(S,T,V)$, and put
\[
 D=t-s^2/2,\quad B=v-st/2-s^3/12,\quad
 F_0=t^2/2+s^2t/4,\quad L=F_0+B^2/D.
\]
Write $D_0=D(X)>0$, $B_0=B(X)$, $u=s-S$, $w=t-T$ and
$R_L=L(x)-L(X)-\nabla L(X)^\top(x-X)$. Direct cancellation gives
\begin{align}
 DD_0^2R_L
 &= (D_0B-B_0D)^2
 +\frac D2\left[D_0(w+Su/2+u^2/6)-B_0u\right]^2\notag\\
 &\quad+\frac{(DD_0u)^2}{12}
 +\frac{DD_0(D_0u)^2}{6}
 +\frac{D(D_0u^2)^2}{36}.\label{eq:rational-remainder}
\end{align}
This is a q-module certificate of the rational first-order type developed
in \cite{nierational}, with denominator multipliers $D(x)$ and $D(X)^2$.
The concave domain generator $D$ has negative Taylor remainder $u^2/2$.

For a direct construction, let $Z$ be a linear functional on polynomials
in $(s,t,v)$ of degree at most six, and let $V_j$ be the monomial vector
through degree $j$. Consider
\begin{align*}
 &Z(V_3V_3^\top)\succeq0,\quad Z(DV_2V_2^\top)\succeq0,\quad Z(D)=1,\\
 &(m,z,q)=(Z(Ds),Z(Dt),Z(Dv)),\quad r\ge Z(DF_0+B^2).
\end{align*}
There are 84 raw functional values and Gram blocks of orders 20 and 10.
Its projection is precisely $D(m,z)>0$, $r\ge L(m,z,q)$.
Indeed, writing $a=Z(Ds^2)$, the second Gram block gives $a\ge m^2$.
The first, restricted to $(1,D)$, gives
\[
 \begin{pmatrix}Z(1)&1\\1&Z(D^2)\end{pmatrix}\succeq0,
 \qquad D(m,z)=Z(D^2)+(a-m^2)/2>0.
\]
Apply $Z$ to \eqref{eq:rational-remainder} with the reference equal to
the output coordinates. The affine Taylor terms vanish after defining $\Lambda(h)=Z(Dh)$ for
polynomial $h$ with $\deg(Dh)\le6$. Here the notation
$\Lambda(L):=Z(DF_0+B^2)$ is defined by the polynomial $DL$;
it does not extend $Z$ to arbitrary rational functions. The two Gram
blocks make every remaining term nonnegative, yielding
$r\ge\Lambda(L)\ge L(m,z,q)$.
Conversely, $Z(P)=P(m,z,q)/D(m,z)$ supplies feasible rank-one Gram blocks
for every point of that strict-domain epigraph.

The strict inequality matters: this rational lift has no finite auxiliary
solution at $D=0$. Its closure is the full set $G(y)\succeq0$; at $D=0$
the added points satisfy $B=0,r\ge F_0$ and are limits with $D>0,B=0$.
The established closure theorem for spectrahedral shadows
\cite[Corollary 3.4, author version]{gn} therefore gives a finite SDP
existence route for this closed set and its intersection with the affine
complementary copy, namely $\K_4$. Closure is not being identified with
the original rational projection, and this argument does not supply the
seven-block bound by itself. We retain the direct polynomial proof in
Theorem~\ref{thm:lift} for that explicit size and boundary statement.

\subsection{Complete fifth-order lift}\label{sec:five}
Write $y=(m,z,q,r,v)$, $x=(m,z,q,r)$ and
$\bar y=(1,\tfrac12,\tfrac13,\tfrac14,\tfrac15)-y$. In this section put
\begin{align}
 \Delta&=mq-z^2-m^4/12,& N&=mr-zq-m^3z/6,& p&=m\Delta,\notag\\
 H(x)&=q^2/m+m^2q/12+mz^2/4-m^5/720,&
 F(x)&=H(x)+N^2/(m\Delta).\label{eq:five-boundary}
\end{align}
The product $pF=pH+N^2$ is polynomial: multiplication by $p=m\Delta$
clears the denominators in $F$. The functionals below act on this polynomial.

The distinction from a direct use of rational-epigraph results is
checkable. The automatic univariate case in \cite[Theorem 3.6, 2009 author
revision]{nierational} does not apply to a function of four prefix moments.
The original integration variable being scalar does not make those moments
one decision variable. The general multivariate framework requires
domain-weighted first-order certificates and additional domain and boundary
hypotheses; convexity of the moment body alone does not verify them.
Nor can the determinant used below be inserted as
a concave polynomial domain constraint: the $(m,q)$ Hessian block of
$\Delta=mq-z^2-m^4/12$ is
$\left(\begin{smallmatrix}-m^2&1\\1&0\end{smallmatrix}\right)$, with determinant
$-1$. This excludes that direct substitution, not all possible applications
of established methods. The functional construction below is proved
directly, including finite boundary witnesses.

\begin{lemma}[Classical interior fifth-moment boundary]\label{lem:five-boundary}
For $x\in\operatorname{int}\K_4$, both $p(x),p(\bar x)>0$, and
$(x,v)\in\K_5$ exactly when $F(x)\le v\le1/5-F(\bar x)$.
\end{lemma}
\begin{proof}
Define $f(x)=\min\{v:(x,v)\in\K_5\}$. Compactness of $\K_5$ and its
projection onto $\K_4$ give attainment on every feasible fiber. Convexity
of $\K_5$ makes $f$ convex and finite on $\operatorname{int}\K_4$;
therefore a subgradient $\lambda$ exists at each such prefix. Its supporting
inequality makes a minimizing density a global minimizer of
$\int_0^1(t^4-\sum_{j=0}^3\lambda_jt^j)\rho(t)\,dt$ over $0\le\rho\le1$.
It equals the quartic's negative-set indicator almost everywhere.
If this density had $k\le3$ genuine interior switches $a_i$, choose the
sign of $W(t)=\pm\prod_{i=1}^k(t-a_i)$ to be positive exactly where the
density is one. For $k=0$, take the appropriate nonzero constant.
The nonzero coefficient vector of $W$ would support $\K_4$ at $x$,
contradicting interiority. A quartic has at most four switches, so it has
four distinct simple roots in $(0,1)$. Since it is monic, its negative set
consists of two strictly interior, disjoint intervals.

For intervals $[a_i,b_i]$, the classical exponential transform is
\[
 1-\exp\left(-\sum_{j\ge0}y_jw^{j+1}\right)
 =1-\prod_i\frac{1-b_iw}{1-a_iw}
 =w\sum_i\frac{c_i}{1-a_iw},\qquad c_i>0.
\]
For $a_1<b_1<a_2<b_2$, the residues are explicitly
\[
 c_1=\frac{(b_1-a_1)(b_2-a_1)}{a_2-a_1}>0,\qquad
 c_2=\frac{(a_2-b_1)(b_2-a_2)}{a_2-a_1}>0.
\]
This is standard Markov-moment algebra \cite{gr}. Consequently the three-by-three Hankel
matrix of the first five transformed coefficients has rank two, with
positive definite leading two-by-two block. Its congruence by
\[
 T=\begin{pmatrix}1&0&0\\m/2&1&0\\z/2+m^2/12&m/2&1\end{pmatrix}
\]
is
\[
 \begin{pmatrix}
 m&z&q\\
 z&q-m^3/12&r-m^2z/6\\
 q&r-m^2z/6&v-m^2q/12-mz^2/4+m^5/720
 \end{pmatrix}.
\]
Eliminating the first row and column gives entries $\Delta/m$, $N/m$ and
$v-H$, so the rank equality gives $v=F(x)$. The leading determinant is
$\Delta>0$. Complementation gives the upper extremum and its positive
denominator. Convex combinations of the two extremizers fill the fiber.
\end{proof}

For a reference $u=(M,Z,Q,R)$ put $h=m-M$, $\Delta_0=\Delta(u)$,
$N_0=N(u)$ and
\begin{align*}
 A_*&=\Delta_0N-N_0\Delta,\\
 B_*&=\Delta_0(Mq-mQ+mM^2h/12)-N_0(Mz-mZ),\\
 C_*&=\Delta_0(Mz+(m-2M)Z)-N_0Mh,\\
 D_*&=\Delta_0((2m-M)z-mZ)-N_0mh.
\end{align*}
Writing $R_F(x,u)=F(x)-F(u)-\nabla F(u)^\top(x-u)$, the joint identity is
\begingroup\interdisplaylinepenalty=10000
\begin{align}
 p(x)p(u)^2R_F(x,u)
 &=(MA_*)^2+\Delta B_*^2+\frac{pM}{6}C_*^2
       +\frac{\Delta}{12}(MD_*)^2\notag\\
 &\quad+\frac{(\Delta M\Delta_0h)^2}{12}
       +\frac{p(M\Delta_0)}6(\Delta_0h)^2\notag\\
 &\quad+\frac{pM}{120}(M\Delta_0h^2)^2
       +\frac{\Delta}{180}(mM\Delta_0h^2)^2.\label{eq:five-certificate}
\end{align}
\endgroup
Coefficient expansion proves this identity with independent $x,u$. It is a
specific domain-weighted first-order certificate in the established rational
lifting framework \cite{nierational}, not a claim to invent that framework.

For a linear functional $\ell$, impose the following seven PSD matrices.
Each row denotes $\ell(gbb^\top)\succeq0$.
\begin{center}
\begin{tabular}{clc}
\toprule
Weight $g$ & Polynomial feature vector $b^\top$ & Order\\
\midrule
$1$ & $(N,\Delta)$ & 2\\
$1$ & $(p,\Delta)$ & 2\\
$\Delta$ & $(q,m,m^2,z)$ & 4\\
$\Delta$ & $(mz,z,m,m^2)$ & 4\\
$\Delta$ & $(m^3,m^2,m)$ & 3\\
$p$ & $(z,m,1)$ & 3\\
$p$ & $(m^2,m,1)$ & 3\\
\bottomrule
\end{tabular}
\end{center}
Use a common functional on the span of all these entries and
$(p,pm,pz,pq,pr,pF)$. This polynomial space has the following basis of size 19:
\begin{align*}
 &(p,pm,pz,pq,pr,pF,\ N^2,N\Delta,\Delta^2,p^2,p\Delta,\Delta q^2,\\
 &\hspace{14mm}\Delta m^2q,\Delta qz,\Delta m^3,\Delta m^4,
           \Delta m^2z,\Delta m^2z^2,pz^2).
\end{align*}
Expansion of monomial coefficients verifies both linear independence and
that every required entry belongs to this span. Thus its functional values
are ordinary finite affine decision variables; the polynomial features are
not substituted nonlinearly into output coordinates. For direct implementation,
Section~\ref{app:coefficients} lists every affine matrix entry in these
19 coordinates; no polynomial elimination or choice of a Gram completion is
left to the reader.

\begin{theorem}[Complete fifth-order lift]\label{thm:five-lift}
A point $y=(x,v)$ belongs to $\K_5$ if and only if its prefix satisfies the
seven-block lift of Theorem~\ref{thm:lift}, and there are two functionals
with the seven PSD matrices above such that
\[
 \ell_-(p)=\ell_+(p)=1,\qquad
 \ell_-(px_i)=x_i,\quad\ell_+(px_i)=\bar x_i,
\]
\[
 \ell_-(pF)\le v,\qquad \ell_+(pF)\le1/5-v.
\]
Each functional acts on its own copy of the formal polynomial variables;
its coordinates are not required to be point evaluations at the output.
The model has 34 auxiliary scalars, 21 PSD blocks of maximum order four,
and two scalar inequalities, including all singular boundary points.
\end{theorem}
\begin{proof}
If the output prefix is interior, set the reference in
\eqref{eq:five-certificate} equal to that output and apply $\ell_-$. The
normalization and output equations cancel its linear Taylor term. All eight
squares lie in the seven displayed spaces and all reference weights are
positive. Hence $\ell_-(pF)\ge F(x)$. The upper copy similarly yields the
upper boundary in Lemma~\ref{lem:five-boundary}.

For a feasible tuple with boundary prefix, mix the entire lifted tuple with
the feasible point-evaluation tuple of the density $1/2$. Affinity preserves
all constraints and places its prefix in $\operatorname{int}\K_4$ for any
positive mixing weight. The preceding argument puts the mixed output in
$\K_5$; compactness then puts the original output there. This proves
sufficiency without assuming that an arbitrary spectrahedral projection is closed.

For necessity let $y$ be physical and mix its density with $1/2$, obtaining
$y_\epsilon$. Its prefix is $x_\epsilon=(1-\epsilon)x+\epsilon c$, where
$c=(1/2,1/4,1/6,1/8)$ is the half-density center. For $\epsilon>0$, the functional
$\ell_\epsilon(f)=f(x_\epsilon)/p(x_\epsilon)$ supplies positive multiples
of rank-one matrices and the required inequalities. We show that its values
on the displayed basis have finite limits, even when raw monomial values diverge.

Extend a density by zero outside $[0,1]$ for the following variance comparison.
If $m>0$ and $\Delta=0$, equality in the bounded-density variance bound
$q-z^2/m\ge m^3/12$ forces a single interval of length $m$, centered at
$z/m$. This follows by comparing the density with that interval under the
weight $(t-z/m)^2-m^2/4$, whose sign agrees with the comparison outside
and inside the interval. Equality forces equality of the densities almost
everywhere. Its third moment gives $N=0$. The function
$\kappa=\Delta/m=q-z^2/m-m^3/12$ is concave for $m>0$:
its $(m,z)$ Hessian has negative first diagonal entry and determinant one,
and its $q$ dependence is affine. At the center
$c=(1/2,1/4,1/6,1/8)$, $\kappa(c)=1/32$. Thus, for sufficiently small
$\epsilon>0$,
\[
 \kappa(x_\epsilon)\ge\epsilon/32,\qquad
 \Delta(x_\epsilon)\ge m\epsilon/64,\qquad
 N(x_\epsilon)=O(\epsilon),\qquad N(x_\epsilon)^2/p(x_\epsilon)=O(\epsilon).
\]
Here $m$ is the fixed boundary mass and $m_\epsilon\ge m/2$.
In the displayed basis order, the normalized coordinates are
\[
 \begin{aligned}
 &(1,m,z,q,r,F,\ N^2/p,N/m,\Delta/m,p,\Delta,q^2/m,\\
 &\hspace{14mm}mq,qz/m,m^2,m^3,mz,mz^2,z^2),
 \end{aligned}
\]
with all entries evaluated at $x_\epsilon$. The only singular quotient
not controlled by the positive limiting mass is $N^2/p$, already bounded
above. Thus all 19 coordinates have finite limits in this case.

If $m=0$, the physical prefix is zero. On its central mixing path,
$m,z,q,r=O(\epsilon)$, $N,\Delta=O(\epsilon^2)$, and $p$ has positive
leading order $\epsilon^3$. The same seven normalized matrices and $F$
are bounded: $m_\epsilon=\epsilon/2$, so $q^2/m$, $qz/m$, $N/m$ and
$\Delta/m$ are $O(\epsilon)$, as are $H$ and $N^2/p$; the other nonconstant
coordinates in the displayed tuple also tend to zero. Every selected
basis element divided by $p(x_\epsilon)$ is therefore a bounded
rational function of $\epsilon$, with a removable singularity. Its limit
defines a finite functional; common polynomial relations persist by linearity.
All PSD constraints and output inequalities persist as finite limits. The
complementary copy is identical, and the prefix uses its finite atomic
witness. This proves necessity at every boundary, not just closure existence.

For clarity, the finite single-interval witness for $m>0,\Delta=0$,
in the displayed 19-coordinate order, is
\[
 \begin{aligned}
 \theta^*={}&(1,m,z,q,r,H,0,0,0,0,0,\ q^2/m,mq,qz/m,\\
             &\hspace{24mm}m^2,m^3,mz,mz^2,z^2).
 \end{aligned}
\]
The physical fifth moment is $v=H$; the first two blocks vanish and the
last five are the weighted outer products above. This includes full density
by setting $(m,z,q,r)=(1,1/2,1/3,1/4)$. At zero density the central path has
$\Delta_\epsilon=N_\epsilon=\epsilon^2(4-\epsilon^2)/192$ and
$p_\epsilon=\epsilon^3(4-\epsilon^2)/384$; its coordinate limit is
$(1,0,\ldots,0)$. The last two blocks are then $\operatorname{diag}(0,0,1)$
and the other blocks vanish. If $m>0,\Delta>0$, point evaluation itself
has no singular denominator, even when the prefix is on another boundary.

The first five values of each 19-dimensional functional are fixed by its
normalization and prefix outputs. Hence there are $14+14+6=34$ auxiliaries.
The two copies of seven blocks and the seven prefix blocks give the stated size.
\end{proof}

\begin{corollary}[A finite-SDP/non-SOCP separation]\label{cor:five-separation}
The complete fifth-order body satisfies $3\le\sxdeg(\K_5)\le4$.
It has a finite affine SDP lift but no lift over any finite product of Lorentz
cones. The same is true of its affine-equivalent integrator reach set.
\end{corollary}
\begin{proof}
Combine Theorem~\ref{thm:five-lift} with Theorem~\ref{thm:lower} below and
the invertible map \eqref{eq:affine}. The exact minimum, three or four,
is not determined here.
\end{proof}

\subsubsection{Explicit coefficient interface}\label{app:coefficients}
Index the 19 basis polynomials in Section~\ref{sec:five} by $b_0,\ldots,b_{18}$
in their displayed order and write $\theta_i=\ell(b_i)$. These are independent
functional coordinates, not evaluations of $b_i$ at the output point.
Define the four linear abbreviations
\begin{align*}
\gamma_{1}&=\theta_{3}-\theta_{8}-\theta_{15}/12,\\
\gamma_{2}&=6(\theta_{4}-\theta_{13}-\theta_{7}),\\
\gamma_{3}&=12(\theta_{12}-\theta_{10}-\theta_{18}),\\
\gamma_{4}&=180(\theta_{5}-\theta_{6}-\theta_{11})-15\theta_{9}-60\theta_{17}.
\end{align*}
The seven pencils, in the feature-table order, are
\begin{align*}
M_{1}(\theta)&=\begin{pmatrix}
\theta_{6}&\theta_{7}\\\theta_{7}&\theta_{8}
\end{pmatrix},&
M_{2}(\theta)&=\begin{pmatrix}
\theta_{9}&\theta_{10}\\\theta_{10}&\theta_{8}
\end{pmatrix},\\[2mm]
M_{3}(\theta)&=\begin{pmatrix}
\theta_{11}&\theta_{3}&\theta_{12}&\theta_{13}\\
\theta_{3}&\theta_{1}&\theta_{14}&\theta_{2}\\
\theta_{12}&\theta_{14}&\theta_{15}&\theta_{16}\\
\theta_{13}&\theta_{2}&\theta_{16}&\gamma_{1}
\end{pmatrix},&
M_{4}(\theta)&=\begin{pmatrix}
\theta_{17}&\theta_{18}&\theta_{16}&\gamma_{2}\\
\theta_{18}&\gamma_{1}&\theta_{2}&\theta_{16}\\
\theta_{16}&\theta_{2}&\theta_{1}&\theta_{14}\\
\gamma_{2}&\theta_{16}&\theta_{14}&\theta_{15}
\end{pmatrix},\\[2mm]
M_{5}(\theta)&=\begin{pmatrix}
\gamma_{4}&\gamma_{3}&\theta_{15}\\
\gamma_{3}&\theta_{15}&\theta_{14}\\
\theta_{15}&\theta_{14}&\theta_{1}
\end{pmatrix},&
M_{6}(\theta)&=\begin{pmatrix}
\theta_{18}&\theta_{16}&\theta_{2}\\
\theta_{16}&\theta_{14}&\theta_{1}\\
\theta_{2}&\theta_{1}&\theta_{0}
\end{pmatrix},\\[2mm]
M_{7}(\theta)&=\begin{pmatrix}
\gamma_{3}&\theta_{15}&\theta_{14}\\
\theta_{15}&\theta_{14}&\theta_{1}\\
\theta_{14}&\theta_{1}&\theta_{0}
\end{pmatrix}.
\end{align*}
Require all seven matrices to be PSD. For the lower copy, set
$(\theta_0,\ldots,\theta_4)=(1,m,z,q,r)$ and impose $\theta_5\le v$.
The upper copy fixes these coordinates to
$(1,1-m,1/2-z,1/3-q,1/4-r)$ and imposes $\theta_5\le1/5-v$.
Each copy has its own $\theta_5,\ldots,\theta_{18}$. Add the seven-block
four-moment prefix on $(m,z,q,r)$: this gives the complete 21-block model.
For synthesis replace these fixed coordinates by the homogeneous values in
Section~\ref{sec:five-synthesis}; the seven displayed pencils do not change.

Substitution of the basis polynomials for the $\theta_i$ recovers each
weighted outer product in the feature table exactly. In particular the
four abbreviations are the polynomial relations for $\Delta z^2$,
$\Delta m^3z$, $\Delta m^5$ and $\Delta m^6$, respectively. This checks
the shared relations, rather than assigning separate variables to repeated
entries. Exact coefficient rank 19 verifies that these values consistently
define a functional on the required span. This is an executable
description of Theorem~\ref{thm:five-lift}, not an additional result.

\subsection{Necessary PSD block size}\label{sec:lower}
Let $P_n([0,1])$ denote the coefficient cone of real polynomials of degree
at most $n$ nonnegative on $[0,1]$.

\begin{lemma}[Normal-cone transfer]\label{lem:transfer}
If a nonempty convex set $C$ has an affine PSD lift with maximum block
order $r$, then for every $v\in C$, its outward normal cone $N_C(v)$ has
such a lift with maximum order at most $\max(r,1)$.
\end{lemma}
\begin{proof}
Write $C=\{LX+c_0: AX=b,\ X\in Q\}$ for a product $Q$ of PSD cones.
Any free scalar variables in the original affine lift can first be encoded
as differences of two nonnegative scalars, adding only order-one blocks.
Compress each block to the span of the ranges of all feasible block matrices.
Finitely many feasible tuples span these spaces; their average is positive
definite on each compressed space, since the kernel of a sum of PSD matrices
is the intersection of their kernels. This produces a strictly feasible
lift with no larger blocks, representing the same set. Zero blocks are
discarded. This is the standard proper-face reduction underlying lift
duality \cite{gpt}.

For $\lambda\in N_C(v)$, maximize $\ip{L^*\lambda}{X}$ over the reduced
lift. The value $\ip{\lambda}{v-c_0}$ is finite and attained. Strict primal
feasibility gives an attained dual with equal value. Hence
\begin{equation}\label{eq:normal-lift}
 N_C(v)=\{\lambda:\exists\eta,
 A^*\eta-L^*\lambda\in Q,\quad
 b^\top\eta=\lambda^\top(v-c_0)\}.
\end{equation}
Conversely the displayed constraints imply
$\lambda^\top(LX+c_0-v)=-\ip{A^*\eta-L^*\lambda}{X}\le0$
for every feasible $X$. All constraints in \eqref{eq:normal-lift} are
affine and use the same PSD block orders. Free variables may be represented
by differences of order-one nonnegative variables.
\end{proof}

\begin{theorem}[Block-size obstruction]\label{thm:lower}
For every $d\ge1$, $\sxdeg(\K_d)\ge\lceil d/2\rceil$.
For $d\ge5$, $\K_d$ has no affine lift over any finite product of Lorentz
cones, including cones of arbitrary finite dimension.
\end{theorem}
\begin{proof}
At the saturated-density point $v_j=1/(j+1)$,
\begin{equation}\label{eq:normal}
 N_{\K_d}(v)=P_{d-1}([0,1]).
\end{equation}
Indeed a nonnegative polynomial $p_\lambda$ gives
$\int p_\lambda(\rho-1)\le0$ for every admissible density. A negative
value gives an interval of negative values; setting $\rho=0$ there and
$\rho=1$ elsewhere violates the normal inequality.

Set $k=\lfloor(d-1)/2\rfloor$. For $k\ge1$, we apply Averkov's polynomial
neighborliness obstruction \cite[Theorem 2, author-version numbering]{averkov}.
The theorem applies to a closed convex cone of bounded-degree polynomials
nonnegative on a domain with nonempty interior.
Suppose there are arbitrarily large finite sets $S$ in that domain.
For every $k$-element subset $A\subset S$, suppose there is a polynomial in
the cone that vanishes on $A$ and is strictly positive on $S\setminus A$.
Then the extension degree is greater than $k$.
Here $[0,1]$ has nonempty interior, and $P_{2k}([0,1])$ is closed and convex
as an intersection of point-evaluation halfspaces. For any finite
$S\subset(0,1)$ and any $k$-element subset $A\subset S$, the polynomial
$\prod_{a\in A}(t-a)^2$ belongs to this cone, is zero on $A$ and is strictly
positive on $S\setminus A$. Such sets $S$ can have arbitrarily large cardinality,
so $\sxdeg(P_{2k}([0,1]))\ge k+1$. If $d-1=2k+1$, setting the highest coefficient
to zero makes $P_{2k}$ a linear slice of $P_{d-1}$; otherwise it is the
same cone. Equation~\eqref{eq:normal} and Lemma~\ref{lem:transfer} yield
$\sxdeg(\K_d)\ge k+1$. The cases $d=1,2$ use the definitional lower bound
one. The polynomial obstruction, rather than any finite numerical test,
provides the universal quantifier in this argument.

Every finite-dimensional Lorentz cone has a lift using finitely many
three-dimensional Lorentz cones, by introducing successive partial
Euclidean norms. The three-dimensional Lorentz cone is isomorphic to $\PSD^2$ by
$(a,b,c)\mapsto\left(\begin{smallmatrix}a+b&c\\c&a-b\end{smallmatrix}\right)$.
An SOC lift would therefore imply $\sxdeg(\K_d)\le2$, impossible for
$d\ge5$.
\end{proof}

Combined with Theorem~\ref{thm:lift}, this gives
$2\le\sxdeg(\K_4)\le3$, not an exact fourth-order classification.
\begin{remark}[Normal-cone limitation]\label{rem:normal-limit}
Changing the boundary point cannot strengthen the normal-cone argument to
order three for $\K_4$. In fact, let $y\in\partial\K_d$ and let
$a_1<\cdots<a_k$ be the genuine interior switches of its representing
density. Every nonzero supporting polynomial has a unique maximizing
density, namely its positive-sign indicator, since its zero set has measure
zero. Choose $W(t)=\epsilon\prod_i(t-a_i)$ with that sign pattern. Then
\[
 N_{\K_d}(y)=\{\operatorname{coeff}(Wq):q\in P_{d-1-k}([0,1])\}.
\]
Indeed a normal polynomial has the required weak signs on every constant
density interval, must vanish at each switch, and its quotient by $W$ is
nonnegative by continuity. The converse follows by pointwise maximization.
Multiplication by the fixed $W$ is a linear isomorphism onto this cone's span.
For $d=4$, all residual degrees are at most three. The classical interval
representation
\[
 q(t)=t(1,t)Q_0(1,t)^\top+(1-t)(1,t)Q_1(1,t)^\top,
 \qquad Q_0,Q_1\in\PSD^2,
\]
represents $P_3([0,1])$; lower degrees are affine coefficient slices. Hence
every ordinary normal cone of $\K_4$ is SOC representable. This is a limit
of this proof method, not a construction of an SOC lift for $\K_4$.
Conversely, the isolated homogeneous Hankel block $H$ is a known non-SOC
slice \cite[Theorem 5]{fawzi},
but no implication from that auxiliary cone to the
projected body is available. Neither observation settles the gap.
\end{remark}

By \eqref{eq:affine}, the obstruction applies to full-state integrator
reach sets with nondegenerate input intervals, not every partial observation
or conditional slice.

\subsection{Exact minimum-peak formulation}\label{sec:synthesis}
\subsubsection{Fourth-order prefix model}
First specialize the problem of Section~\ref{sec:control-problem} to
fourth-order chains. This model supplies the homogeneous prefix used by
the fifth-order construction; all times and physical rows remain fixed.

Homogenize \eqref{eq:H}--\eqref{eq:Q} by replacing the upper-left constant
one in $H,J_+,J_-$ with $J$, leaving all other entries affine in scaled
moments $Y$ and auxiliaries $Z$. Denote these pencils by $\widehat M_k$.
\begin{lemma}\label{lem:homogeneous}
\[
 \{(J,Y):J\ge0,\ \exists Z,\ \widehat M_k(J,Y,Z)\succeq0\ \forall k\}
 =\{(J,Jy):J\ge0,\ y\in\K_4\}.
\]
\end{lemma}
\begin{proof}
For $J>0$, divide every pencil by $J$ and apply Theorem~\ref{thm:lift}.
At $J=0$, the zero first diagonal of $H$ forces $m=a=0$, and its second
row then forces $b=0$, while $c\ge0$. The zero first rows of $J_\pm$
force $z=0$; $T_-$ then forces $d=0$. The last diagonals of $T_-$ and
$T_+$ give $e-c/2\ge0$ and $-e-c/2\ge0$, so $c=e=0$.
The last diagonal of $J_\pm$ gives $f\ge0$. Matrix $Q_-$ forces $q=0$;
the two $Q$ last diagonals give $r-f/2\ge0$ and $-r-f/2\ge0$, so
$f=r=0$. There are no nonzero apex directions. The zero tuple is feasible.
\end{proof}

Partition $[0,T]$ at the row times, obtaining $N$ segments $[l,r]$.
On a segment of length $h$, set
$Y_j=\int_0^1 s^j[u_a(l+hs)+J]/2\,ds$.
Its signed input moments are $2Y_j-J/(j+1)$. Remove the known zero-input
flow from each row and write the resulting target as $\beta_i$.
The convolution kernel for derivative index $d=0,1,2,3$ is
\[
 k_{ia}(t)=\mathbf1_{t\le t_i}\sum_{d=0}^3
       w_{i,a,d}\frac{(t_i-t)^{3-d}}{(3-d)!}.
\]
Expand $h k_{ia}(l+hs)=\sum_{j=0}^3c_{iaj}^{[l,r]}s^j$ on each segment.
The row constraints are exactly
\begin{equation}\label{eq:finite-rows}
 \left|\sum_{a,[l,r],j}c_{iaj}^{[l,r]}
       \left(2Y_{aj}^{[l,r]}-\frac{J}{j+1}\right)-\beta_i\right|
       \le\epsilon_i.
\end{equation}

\begin{corollary}\label{cor:synthesis}
Minimizing $J$ subject to $J\ge0$, the homogenized seven-block constraints
on every axis/segment, and \eqref{eq:finite-rows}, is equivalent to the
continuous-input problem \eqref{eq:waypoints}. There are $10AN+1$ scalar
variables, $3AN$ PSD blocks of order three and $4AN$ of order two.
\end{corollary}
\begin{proof}
Each feasible input gives feasible moments. Conversely Lemma~\ref{lem:homogeneous}
gives a bounded measurable input on each segment; concatenate these inputs.
Convolution gives a continuous state trajectory satisfying all rows.
The peak-zero case follows from the apex argument. Thus the feasible peak
sets and optimal values coincide. If feasible, a bounded minimizing sequence
has a weak-star convergent subsequence and the finitely many kernel integrals
pass to its limit, so the minimum is attained. No input mesh or prescribed
switch count is introduced in the outer model.
\end{proof}

\subsubsection{Fifth-order model and the zero-peak apex}\label{sec:five-synthesis}
Replace the fourth-order dynamics above by fifth-order chains. Homogenize
the complete lift in Theorem~\ref{thm:five-lift}, using functional coordinates
$(J,Y_0,Y_1,Y_2,Y_3,\zeta_-)$ on the lower side and
$(J,J-Y_0,J/2-Y_1,J/3-Y_2,J/4-Y_3,\zeta_+)$ on the upper side,
where each $\zeta$ has 14 entries. The prefix is the homogeneous four-moment
lift; the last two scalar rows are
$Y_4\ge\ell_-(pF)$ and $J/5-Y_4\ge\ell_+(pF)$.
All entries are affine decision expressions.

\begin{lemma}[Exact fifth-order homogenization]\label{lem:five-homogeneous}
The output projection of these constraints with $J\ge0$ is precisely
$\{(J,Jy):J\ge0,\ y\in\K_5\}$.
\end{lemma}
\begin{proof}
For $J>0$ divide by $J$ and use Theorem~\ref{thm:five-lift}.
At $J=0$, Lemma~\ref{lem:homogeneous} forces $Y_0=\cdots=Y_3=0$.
The polynomial identity
\begin{equation}\label{eq:five-apex}
 pF=N^2+\Delta q^2+\frac{p^2}{12}
       +\frac{\Delta(mz)^2}{3}+\frac{\Delta m^6}{180}
\end{equation}
expresses its functional value as a nonnegative combination of retained PSD
diagonal entries. Thus both $\ell_\pm(pF)$ are nonnegative, even at the apex.
The two scalar rows force $Y_4=0$. The all-zero lifted tuple is feasible;
uniqueness of the auxiliary values is unnecessary.
\end{proof}

For fifth-order dynamics the physical kernel is
\[
 k_{ia}(t)=\mathbf1_{t\le t_i}\sum_{d=0}^4
 w_{i,a,d}\frac{(t_i-t)^{4-d}}{(4-d)!}.
\]
Its local polynomial expansion uses five signed moments per axis/segment.
Equation~\eqref{eq:finite-rows} therefore applies with $j=0,\ldots,4$.
\begin{corollary}\label{cor:five-synthesis}
The fifth-order continuous minimum-peak problem with constraints
\eqref{eq:waypoints} is equivalent to a finite affine SDP with $39AN+1$
variables, $21AN$ PSD blocks of maximum order four, $2AN$ scalar cone rows,
and the physical waypoint rows. The outer formulation has no input-time mesh.
\end{corollary}
\begin{proof}
Use Lemma~\ref{lem:five-homogeneous} on every axis/segment and concatenate
the corresponding measurable inputs. The same convolution and weak-star
compactness argument as for Corollary~\ref{cor:synthesis} gives identical
feasible peak values and attainment for feasible tasks. This proves the
zero-peak case as well; it is not a division-by-zero convention.
\end{proof}

\subsection{Shared certification and continuous-dual reference}\label{sec:certification}
The representation and synthesis theorems allow the real data specified in
Section~\ref{sec:control-problem}. The exact-rational implementation below
instead certifies instances whose horizon, observation times, initial states,
row coefficients, targets and tolerances are specified exactly as rationals.
Its control amplitudes, breakpoints and checked dual multipliers are rational
as well. Replacing real data by finite decimals certifies that rational
instance, not the original real-data equalities; transferring the certificate
would require a separate error enclosure. Rational input data do not guarantee
success of the finite-budget recovery.

Let $k_a(t)=\sum_i\lambda_i k_{ia}(t)$ for any vector $\lambda$ and define
\[
 I(\lambda)=\sum_a\int_0^T|k_a(t)|\,dt,\qquad
 B(\lambda)=\lambda^\top\beta-\sum_i\epsilon_i|\lambda_i|.
\]
Every feasible input has $B(\lambda)\le JI(\lambda)$. Hence, if
$\widehat I\ge I$ and $\widehat I>0$,
\begin{equation}\label{eq:dualbound}
 L=\max\{0,B(\lambda)/\widehat I\}\le J_*.
\end{equation}
If $I(\lambda)=0$, feasibility requires $B(\lambda)\le0$; such a vector
provides no positive lower bound through \eqref{eq:dualbound}. A positive $B$ with zero $I$
would instead certify infeasibility. The zero lower bound is always available
for a feasible problem.
The usual norm dual is $\sup\{B(\lambda):I(\lambda)\le1\}$.
For completeness, let $Z$ be the compact image of the unit input ball in
measurement space and $E=\prod_i[-\epsilon_i,\epsilon_i]$. The extended
value function $f(\beta)=\inf\{J\ge0:\beta\in JZ+E\}$ is proper,
convex and lower semicontinuous: its sublevels are the compact sets $cZ+E$.
Its conjugate is $h_E(\lambda)$ when $h_Z(\lambda)\le1$ and infinity
otherwise. Since $h_Z=I$, biconjugacy proves equality of the dual supremum
and primal value. Dual attainment or numerical convergence is not assumed.

The reference implementation integrates absolute cubic or quartic polynomials using
their real roots and analytic antiderivatives, and solves this finite
continuous dual numerically. It is not a coarse input-grid baseline.
Both routes use identical postprocessing. On a rational subdivision of a
segment, the degree-$n$ Bernstein coefficients $b_0,\ldots,b_n$, for
$n=3$ or $4$, give
\[
 \int_l^r|p(t)|\,dt\le\frac{r-l}{n+1}\sum_{j=0}^n|b_j|.
\]
This supplies $\widehat I$ with rational arithmetic. An independent checker
reconstructs kernels and Bernstein coefficients from endpoint values and
derivatives and verifies complete interval coverage. For example, the quartic
middle coefficient is $b_2=p(l)+(r-l)p'(l)/2+(r-l)^2p''(l)/12$;
this independent endpoint formula is not the producer's power-basis conversion.

For a primal upper bound, an inner LP proposes piecewise-constant controls.
Rational correction satisfies selected active rows exactly; all rows are
then checked by exact convolution, and the actual input peak $U$ is
recomputed. No waypoint tolerance is enlarged. The resulting certificate is
$L\le J_*\le U$. The earlier fourth-order suites use the target
\begin{equation}\label{eq:tolerance}
 U-L\le10^{-6}\max(1,U).
\end{equation}
The recovery is not a universal completeness theorem. It can reject a
floating solution or stop before the requested gap is achieved. Such outcomes
do not prove the continuous problem infeasible.

\section{Numerical validation of the theoretical results}\label{sec:validation}
The tests address three consequences of Section~\ref{sec:theory}:
(i) the lift reproduces support values and admits exact boundary witnesses;
(ii) homogeneous synthesis gives checkable bounds for fixed-time state rows;
and (iii) block compression can be compared with another exact formulation.
Recovery ablations and feedback replays identify implementation limits.
Finite simulations do not prove the non-SOC theorem, replace boundary proofs,
or turn model equivalence into a real-time guarantee.

\subsection{Support functions and boundary witnesses}\label{sec:support-validation}
The fixed-body model used for this support-function audit has 39 variables
(five output coordinates and 34 auxiliaries), 21 PSD blocks, and two scalar
inequalities. The synthesis model of Corollary~\ref{cor:five-synthesis}
adds a shared peak variable, giving $39AN+1$ variables, or 40 for one axis
and one segment. An exact coefficient implementation checks finite rational
witnesses at zero, full, single-interval and multiple-switch densities, as well
as all exact support maximizers in the following audit.

Before solving, nine root families were fixed, each with both leading signs:
the empty family; $(1/2)$; $(1/4,3/4)$; $(1/5,1/2,4/5)$;
$(1/5,2/5,3/5,4/5)$; $(1/100,1/3,2/3,99/100)$; $(1/2,1/2)$;
$(1/4,1/4,3/4,3/4)$; and $(49/100,51/100)$. Their polynomials give 18
support directions. The empty family gives the constant polynomials $\pm1$.
All directions are normalized by coefficient Euclidean norm. Exact rational
integration of each positive part gives an independent optimum and maximizing
density. A floating-point polynomial-root oracle is evaluated separately.

One non-warm CLARABEL solve per direction used absolute/relative gap and
feasibility tolerances $10^{-10}$, with at most 300 iterations and one thread.
No retry or solver switching was used. All 18 cases met the predeclared
$10^{-7}$ support-error and PSD/scalar-residual thresholds; the largest
absolute normalized support error was $2.099\times10^{-8}$ (rounded upward).
However, 16 cases returned precision warnings. These numerical checks do not
certify that each floating-point optimizer output is an exactly feasible
physical moment vector. The rational maximizing witnesses are separate
evidence. This support audit alone gives no feedback or real-time guarantee;
Section~\ref{sec:five-feedback} separately examines fifth-order feedback.
The median SDP solve-call time was 3.646\,ms, versus 0.0984\,ms for the
numerical root oracle. One-time model construction and canonicalization took
0.2600\,s and 0.07228\,s; subsequent solve-call times reuse the compiled model.
These measurements are not cold-start control-cycle costs and do not show
an SDP speed advantage for a single support query.

\subsubsection{Singular-limit and membership diagnostics}\label{sec:boundary-paths}
A supplementary study with tasks and tolerances specified before solving
examines the finite-witness argument directly. Six endpoint densities are the indicators
of the empty set, $[0,1]$, $[0,2/5]$, $[3/5,1]$, $[1/4,3/4]$, and
$[0,1/4]\cup[1/2,3/4]$. Each is mixed with density $1/2$ at
$\epsilon\in\{1,10^{-1},10^{-2},10^{-4},10^{-6},10^{-8},10^{-10},0\}$.
For all 48 points, rational evaluation of the 19 selected coordinates on
both sides, with removable limits at zero, satisfies the affine PSD
constraints in Sections~\ref{sec:four} and~\ref{app:coefficients}
and all remaining scalar and prefix constraints of
Theorem~\ref{thm:five-lift} exactly.
The largest selected-coordinate infinity norm is one. At $\epsilon=10^{-10}$,
the largest distance to the limiting coordinates is below $1.5\times10^{-10}$.
This checks particular finite witnesses, not the general boundary proof.

Separately, each fixed-output model minimizes the infinity norm of its 34
auxiliaries, without a supplied witness or warm start. The CLARABEL settings
are those of the preceding support audit. All 48 solves meet the fixed
$10^{-7}$ output, PSD and scalar-residual thresholds, but all return
\texttt{optimal\_inaccurate}. Numerical acceptance is not exact feasibility.
To expose that distinction, use the moment vectors $y$ for the zero density,
the full density and the indicator of $[1/4,3/4]$. Their respective supporting
polynomials are $-1$, $1$ and $-(t-1/4)(t-3/4)$. For each polynomial's
coefficient vector $c$ and
$\eta\in\{10^{-3},10^{-6},10^{-9}\}$, compare the physical point
$(1-\eta)y+\eta(v/2)$ with $y+\eta c/(c^\top c)$.
The outside point $y+\eta c/(c^\top c)$ violates the original body's
support inequality by exactly $\eta$.
All nine inside points pass the numerical diagnostic. All six outside points
at the two larger offsets are reported infeasible, whereas all three outside
points at $10^{-9}$ are numerically accepted with precision warnings.
These are tolerance-level false acceptances, not exact counterexamples to
Theorem~\ref{thm:five-lift}. They explain why support comparisons and floating
residuals cannot replace finite-witness and sufficiency proofs.


\subsection{Fifth-order minimum-peak synthesis}\label{sec:five-pilot}
Fifteen deterministic tasks were specified before any candidate solve: two
zero-input ballistic tasks (one and two axes); integral anchors of peak
$10^{-4}$ and $10^3$; three exact terminal-state anchors whose optimal inputs
have four switches; and eight waypoint tasks from one/two axes, two fixed
seeds, and intermediate tolerances zero/$1/200$. The waypoint tasks specify position
and velocity rows at quarter times and full five-state terminal rows.
Two-axis intermediate rows couple the axes. Anchor 0 uses switches
$(1/5,2/5,3/5,4/5)$ at horizon one and peak one; anchor 1 uses the same
switches at horizon one half and peak $3/2$. Anchor 2 uses switches
$(1/100,1/3,2/3,99/100)$ at unit horizon and peak one.
Known optimum values and generating inputs are used only in subsequent checks,
not as inputs to either proposal or recovery algorithm. Scaled copies are
not independent statistical samples.

Each method makes one fresh candidate solve per task; the first method
alternates by task. CLARABEL uses $10^{-10}$ gap/feasibility tolerances,
300 iterations and one thread; the continuous quartic-root dual uses SLSQP,
600 iterations, analytic gradients and $10^{-11}$ stopping tolerance.
There is no cross-method warm start, retry or outcome-dependent case removal.
Both feed only multipliers to the same recovery, limited to four LP iterations
and 257 mesh nodes. SDP primal moments are not directly converted to inputs.
Every final input must satisfy every original row in exact rational arithmetic.
Initial dual bounds $L_0$ are saved separately because recovery LP multipliers
can improve the final lower bound $L$.

To avoid a fixed absolute tolerance hiding small-input errors, let $U_i$ be
the one-interval quartic Bernstein upper bound on
$\sum_a\int_0^{t_i}|k_{ia}(t)|\,dt$. Define the input-only scale
\[
 S=\max_{i:U_i>0}\frac{(|\beta_i|-\epsilon_i)_+}{U_i},
 \qquad U-L\le10^{-5}S.
\]
When all numerators vanish, zero input is feasible and set $S=1$.
For a feasible nonzero task $0<S\le J_*$, so the criterion is at least
as strict as a $10^{-5}$ relative-to-optimum gap. The checker reconstructs
$S$ from the task, without reading a known optimum or a proposed input.

\paragraph{Results.}
All 30 runs yield valid rational brackets, but only 14/15 SDP and 12/15
continuous-dual runs satisfy the gap target. Table~\ref{tab:five-pilot}
retains every task. A valid but wide bracket is not called converged.
\begin{table}[htbp]
\centering\small
\caption{Every fifth-order task; target gap/$S\le10^{-5}$. Times include
cold candidate construction, optimization, common recovery and final checking.
Values are rounded only for display.}\label{tab:five-pilot}
\begin{tabular}{lrrrr}
\toprule
Task & SDP gap/$S$ & Dual gap/$S$ & SDP (s) & Dual (s)\\
\midrule
Zero, one axis & 0 & 0 & 0.4800 & 0.0252\\
Zero, two axes & 0 & 0 & 0.4987 & 0.0947\\
Integral anchor, $J_*=10^{-4}$ & 0 & 0 & 0.0735 & 0.0059\\
Integral anchor, $J_*=10^3$ & 0 & 0 & 0.1032 & 0.0066\\
Terminal anchor 0 & $3.006\cdot10^{-6}$ & $1.062\cdot10^{-7}$ & 0.2855 & 0.2168\\
Terminal anchor 1 & $5.696\cdot10^{-6}$ & $5.210\cdot10^{-9}$ & 0.2826 & 0.2222\\
Terminal anchor 2 & $4.356\cdot10^{-6}$ & $1.123\cdot10^{-6}$ & 0.4229 & 0.2271\\
1 axis, seed 1, exact & $2.574\cdot10^{-6}$ & $1.814\cdot10^{-3}$ & 0.8325 & 1.2885\\
1 axis, seed 1, tolerance & $2.382\cdot10^{-7}$ & $1.763\cdot10^{-8}$ & 0.6951 & 0.3394\\
1 axis, seed 2, exact & $1.440\cdot10^{-6}$ & $6.571\cdot10^{-7}$ & 0.9616 & 0.7040\\
1 axis, seed 2, tolerance & $4.954\cdot10^{-6}$ & $3.277\cdot10^{-9}$ & 0.8353 & 0.4197\\
2 axes, seed 1, exact & $4.106\cdot10^{-5}$ & $8.759\cdot10^{-4}$ & 4.2392 & 3.6691\\
2 axes, seed 1, tolerance & $1.392\cdot10^{-7}$ & $4.016\cdot10^{-8}$ & 1.7162 & 1.0032\\
2 axes, seed 2, exact & $6.239\cdot10^{-7}$ & $3.465\cdot10^{-4}$ & 1.8710 & 3.9292\\
2 axes, seed 2, tolerance & $1.170\cdot10^{-7}$ & $2.985\cdot10^{-7}$ & 1.7066 & 0.8731\\
\bottomrule
\end{tabular}
\end{table}

Thirteen SDP runs have precision warnings. The dual has three non-success
statuses: line-search failure on the $10^3$ anchor and iteration limits on
both two-axis exact-waypoint tasks. The $10^3$ anchor nevertheless has an exact
optimal certificate. Conversely, the one-axis seed-1 exact task reports
dual solver success but fails the certified gap target. Solver status,
certificate validity and requested accuracy are distinct quantities.
The largest scaled gaps are $4.107\cdot10^{-5}$ and $1.814\cdot10^{-3}$,
respectively (rounded upward). Recovery uses 19 and 21 LP calls, giving 70
total optimization calls. All LP option-forwarding notices remain archived.

In all 14 SDP successes, the initial bound $L_0$ and final upper bound already
meet the target; recovery does not improve the lower bound. The dual's three
failed gaps improve during recovery without closing. These observations test
the implemented multiplier-to-certificate pipeline, not a universal recovery
guarantee.

On the 12 commonly converged tasks, the SDP is slower in every case; the
median paired full-cost ratio is 1.9725. This conditional statistic is given
together with all failures, not as a selected-only comparison. One run per
task does not establish statistical timing or reliability superiority.
Execution timing is examined separately in Section~\ref{sec:five-feedback}.

\subsubsection{Separating the SDP proposal from recovery}\label{sec:five-attribution}
The supplementary study re-solves all 15 tasks once with the unchanged SDP
proposal. Direct binomial convolution of its returned homogeneous moments
checks the original waypoint rows, separately from input recovery. All 15
raw outputs meet the fixed floating diagnostic: row residuals at most
$10^{-7}$ after scaling by the maximum of one, the adjusted target magnitude
and row tolerance, and cone residuals at most $10^{-7}\max(1,|J|)$.
The largest scaled row residual is below $2.185\times10^{-9}$; thirteen
proposals carry precision warnings. These checks do not certify raw moments
as exact physical inputs.

Each proposal is followed by two runs of the same recovery: one seeded by
its multipliers and one starting with zero multipliers, with order alternated
by task. Both keep the original four-LP, 257-node budget and
$10^{-5}S$ accuracy target. The zero-seed route can obtain new multipliers
from its own LPs, so it is not a dual-free algorithm.
Table~\ref{tab:five-attribution} reports all tasks. For all 14 successful
SDP-seeded cases, the initial proposal lower bound $L_0$ and final input
upper bound $U$ already meet the target; recovery does not improve $L_0$.
The remaining SDP-seeded case is the same two-axis exact-waypoint task
that misses the target in Table~\ref{tab:five-pilot}.
\begin{table}[htbp]
\centering\small
\caption{Fresh fifth-order recovery attribution on the same 15 tasks.
An initial-gap pass tests $U-L_0\le10^{-5}S$ using the final recovered $U$;
it is not a pre-recovery primal certificate.}
\label{tab:five-attribution}
\begin{tabular}{lrrrr}
\toprule
Recovery seed & Valid & Initial-gap pass & Final-gap pass & LP calls\\
\midrule
SDP multipliers &15/15&14/15&14/15&19\\
Zero multipliers &15/15&2/15&4/15&48\\
\bottomrule
\end{tabular}
\end{table}

Exact integration of each recovered control reconstructs its segment moments.
For each recovery route, all 53 nonzero-peak segment bodies admit exact
complete-lift witnesses, and both zero-peak tasks have zero homogeneous
moments. This tests the physical-input-to-lift direction and the apex on
actual outputs. The recovered moments need not equal the raw SDP moments:
the largest absolute discrepancy in the SDP-seeded route is about $0.1625$.
Nonunique controls can have different moments, and the recovery does not
enforce their equality. Thus the experiment supports the value of the
SDP-provided multipliers under this fixed budget, not direct reconstruction
of every raw SDP optimizer or SDP-exclusive certification. It comprises
15 fresh synthesis proposals and 67 recovery LPs; together with
Section~\ref{sec:boundary-paths}, the completed supplementary run uses 81 SDP
solves. No new timing superiority or feedback claim is made.

\subsection{Fourth-order formulation-size comparison}\label{sec:matched}
This tests whether the smaller lift in Theorem~\ref{thm:lift} reduces
implementation cost relative to another exact formulation, rather than
against an inexact input grid. Fourteen deterministic problems use the
seven-block model (S), the generic polynomial-coordinate model (G), and
the continuous root dual (D), with three cold repetitions each.
A seeded permutation (20260928) and cyclic method order balance the 126
queries; these are repeated timings, not independent random instances.

The generic model is fully defined in Section~\ref{sec:four}, including its
restriction to the seven-block system and its zero-peak output. The two
exact models therefore share all physical rows, using
$1+10AN$ versus $1+16AN$ variables. At $A=2,N=8$, canonicalization gives
161/257 variables and 517/661 matrix rows, with 36 equalities in both.

The base cases have $A\in\{1,2\}$, $N\in\{2,4,8\}$ and $T=S=1$,
where $S$ is the generating input amplitude. Variants change horizon to
$1/4,4$, amplitude to $10^{-4},10^3$, or use a single switch at
$1-\varepsilon$ with $\varepsilon=10^{-1},\ldots,10^{-4}$ and known optimum
one. Exact intermediate position/velocity rows and terminal full-state rows
are generated from withheld controls; two-axis intermediate rows mix axes.
The full deterministic generation rule and task records are included in
the reproducibility archive described at the end of this section.

The solver and recovery budgets are those of the other synthesis tests.
The scaled acceptance criterion is
\begin{equation}\label{eq:matched-tolerance}
 (U-L)/\max(S,U)\le10^{-6}.
\end{equation}
Recovery receives tolerance $10^{-6}\min(1,S)$ against $\max(1,U)$,
which is no looser than this test. Timings include cold construction,
canonicalization, solution, recovery and independent checking; imports and
one-time symbolic preparation (0.09637 seconds) are excluded.
Raw physical rows are used, with single-thread settings and no competing
benchmark job, rather than optimized warm-start implementations.

All 126 certificates are valid. The target is met on 24/42 S, 24/42 G
and 27/42 D queries, or 8/14, 8/14 and 9/14 problems; each problem's
three repeated outcomes agree. Maximum scaled gaps are
$1.694\cdot10^{-4}$, $1.241\cdot10^{-3}$ and $1.522\cdot10^{-2}$.
S/G have 21/30 inaccurate statuses; D has 27 successful terminations,
six iteration limits and nine line-search failures. There are 417 total
optimizer calls, including 291 recovery LPs. Table~\ref{tab:matched}
retains all problems and failures.
\begin{table}[ht]
\centering\small
\caption{Matched experiment: median full-pipeline time over three repetitions.
Y means all three satisfy \eqref{eq:matched-tolerance}; N means none does.
Times for N entries are unsuccessful pipeline costs, not time to the requested
accuracy. All entries have valid feasible trajectories and bounds.}
\label{tab:matched}
\begin{tabular}{@{}lcrrr@{}}
\toprule
Problem & Target S/G/D & S (s) & G (s) & D (s)\\
\midrule
1 axis, 2 segments & Y/Y/Y & 0.2281 & 0.2450 & 0.1815\\
1 axis, 4 segments & Y/N/Y & 0.7257 & 1.3351 & 0.4363\\
1 axis, 8 segments & Y/Y/Y & 1.1529 & 1.3101 & 1.3755\\
2 axes, 2 segments & Y/Y/Y & 0.4609 & 0.5173 & 0.5645\\
2 axes, 4 segments & Y/Y/N & 1.7214 & 1.7938 & 3.4921\\
2 axes, 8 segments & N/N/N & 7.3479 & 7.6230 & 5.2169\\
$T=1/4$ & Y/Y/N & 1.4623 & 1.5814 & 2.0611\\
$T=4$ & Y/Y/Y & 3.1912 & 2.3725 & 2.1360\\
$\varepsilon=10^{-1}$ & Y/Y/Y & 0.2472 & 0.3480 & 0.0569\\
$\varepsilon=10^{-2}$ & N/N/N & 0.6307 & 0.6072 & 0.4065\\
$\varepsilon=10^{-3}$ & N/N/Y & 0.6733 & 0.7252 & 0.1124\\
$\varepsilon=10^{-4}$ & N/N/Y & 0.6247 & 0.7499 & 0.1182\\
$S=10^{-4}$ & N/N/N & 2.5076 & 2.8292 & 2.4937\\
$S=10^3$ & N/Y/Y & 2.7166 & 2.4570 & 2.4629\\
\bottomrule
\end{tabular}
\end{table}

On the seven jointly successful S/G problems, S is faster in six
(median ratio 0.9247003), but costs 34.5\% more at $T=4$.
On the six jointly successful S/D problems, S is faster in two
(median ratio 1.3753754). Neither comparison establishes general dominance.
S/G median solve-call times are 0.00519/0.00669 seconds, compared with
0.0891/0.1097 for construction, 0.1982/0.2520 for canonicalization and
0.4829/0.6681 for recovery; component medians do not sum to a total median.

All recovery LPs and attempted rational primal corrections succeed.
For SDP saturation failures, lower bounds are within $3\cdot10^{-9}$ of
the known optimum one: recovery upper bounds dominate the remaining gap.
The largest case fails for all methods, with a dual iteration limit and
recovery mesh caps. Small-amplitude failures under raw scaling do not exclude
standard rescaling. These archived diagnostics distinguish model exactness
from finite-budget accuracy without tuned replacement runs.


\subsection{Fourth-order calibration and recovery ablation}\label{app:four-pilot}
This supporting test separates the value of outer-solver multipliers from
the shared recovery procedure. Eighteen predetermined tasks comprise six
scaled rest-to-rest cases, eight one-/two-axis waypoint cases with exact or
tolerant rows, two ballistic zero-peak cases and two total-input anchors
($10^{-4},10^3$). Generating controls are withheld from the methods.
The rest-to-rest optimum $384|d|/T^4$ has the unit-horizon switches
$(1-1/\sqrt2)/2,1/2,(1+1/\sqrt2)/2$ and alternating signs; the terminal
dual vector $(384,-192,40,-4)$ has $\int_0^1|k(t)|\,dt=1$, where $k$
is its aggregate kernel as defined in Section~\ref{sec:certification}.

The SDP and cubic-root dual use the same solver tolerances as the fifth-order
suite. Shared recovery starts with four cells per segment and each method's
own multiplier roots, then refines using LP multipliers and midpoints, with
at most four LP calls and 257 knots. Two ablations instead supply zero
multipliers: one stops after the initial LP and the other uses the full
adaptive budget. Neither receives an optimum or another method's output.
Acceptance uses \eqref{eq:tolerance}; Table~\ref{tab:pilot} reports all four
pipelines.
\begin{table}[ht]
\centering\small
\caption{Independently checked rational certificates. Normalized gap is
$(U-L)/\max(1,U)$. All 18 problems are deterministic and related; the
counts are not population success probabilities.}\label{tab:pilot}
\begin{tabular}{@{}p{50mm}rrr@{}}
\toprule
Pipeline & Valid bounds & Target gap & Max. normalized gap\\
\midrule
SDP plus shared recovery &18/18&18/18&$4.729\cdot10^{-7}$\\
Continuous dual plus recovery &18/18&16/18&$1.868\cdot10^{-4}$\\
Quarter-segment LP only &18/18&4/18&$5.822\cdot10^{-1}$\\
Zero-seed adaptive recovery &18/18&4/18&$7.714\cdot10^{-3}$\\
\bottomrule
\end{tabular}
\end{table}

The two original routes require 36 outer solves and 40 recovery LPs;
the ablations require 18 and 60 LPs. Their four successful tasks are the
zero-peak and total-input anchors. The other 14 adaptive ablation gaps improve
but miss the target. There are 11 inaccurate SDP statuses and, for the dual,
two iteration limits and one line-search failure. The two failed dual gaps
are exact two-axis tasks (about $6.359\cdot10^{-5}$ and
$2.308\cdot10^{-4}$), whose next refinements exceed the knot cap.

For the 16 jointly successful original pairs, SDP is faster in none;
single-run cold medians are 0.1412 and 0.0966 seconds. Fixed execution order
and related cases prevent a statistical speed or population reliability
claim. The ablation supports informative multiplier initialization within
this budget, not SDP-exclusive certification.


\subsection{Feedback execution and timing limits}\label{sec:five-feedback}
\subsubsection{Fifth-order deadline protocol}
Four deterministic scenarios---one-axis nominal, two-axis nominal, two-axis
finite pulse and two-axis alternating disturbance---are each run with both
methods and two execution modes. Every episode has twelve half-second slots,
a terminal-rest horizon of three seconds, a command cap of five, and actuator
quantization with step $10^{-9}$ and ties rounded to even. Initial positions are $1/10$ and $-2/25$;
all other initial derivatives vanish. The pulse is $(1/20,-1/25)$ on slots
3--5. Persistent disturbance amplitudes are $1/30$ and $1/40$, with two positive
then two negative slots and a one-slot phase shift on the second axis.
The planner receives exact current state, not future disturbance values.
Methods alternate order across scenario/mode pairs and retain the same
fixed synthesis and recovery budgets. No missed-deadline trial is rerun
with a larger deadline or different tolerances.

In the ideal-instant reference, a new eligible plan acts immediately, and
plant evolution deliberately ignores computation time. In the one-period
mode, the old quantized queue acts during computation. Its known prefix
and the exact current state predict the next sample's state; the new
terminal-rest problem starts from that prediction. A valid certificate with
peak at most five can replace the queue only at the next boundary and only
if the measured service time is at most $0.5$ seconds. A wider optimality
gap is recorded but does not by itself invalidate feasibility.

A late result is discarded, but its worker remains occupied until its measured
completion time. Sampling requests arriving while busy are skipped, rather
than assuming free cancellation or applying the late result retroactively.
Invalid or over-cap results leave the old queue unchanged. When no plan is
available, or when the old plan is exhausted, the commanded input is zero;
this fallback is not claimed to be safe.
The cap bounds command, not command plus external disturbance.

Service time includes prediction, fresh modeling, optimization, common recovery
and final independent plan checking; it excludes archive serialization and
offline instrumentation. Integer nanoseconds determine deadline comparisons.
These are one-run, one-machine service-time replays, not worst-case execution
time bounds or hard-real-time hardware experiments. Different methods evolve
different states, so their subsequent queries are not identical-input timing
pairs or independent random reliability trials.

An independent rational checker reconstructs queue identities, absolute-time
prefixes, problem data, deadline and busy decisions, inputs and every state
transition. Producer propagation is sequential and checker propagation uses
direct convolution. The local derivative-$j$ discrepancy from the intended
prefix is bounded by $(|d_a|+q/2)\delta^{5-j}/(5-j)!$. For the delayed
prediction, quantization is already included, leaving the bound
$|d_a|\delta^{5-j}/(5-j)!$. These local bounds do not imply stability or
recursive feasibility. In particular, a plan's nominal terminal equality
does not become an equality for the disturbed plant.

\subsubsection{Fifth-order outcomes}
The 16 episodes contain 192 slots, 158 actual planning queries and 410
optimization calls (158 candidates and 252 recovery LPs). All queried plans
have valid rational brackets, but only 137 meet the input-scaled gap target.
Table~\ref{tab:five-feedback-counts} distinguishes accuracy from activation.
There are 56 SDP precision warnings and no discarded warning records.
\begin{table}[htbp]
\centering\small
\caption{Each row spans four scenarios and 48 slots. ``Late'' in ideal mode
is recorded but ignored by plant timing; ``Busy'' counts skipped slots, not
failed solves. Six ideal SDP plans are feasible but wider than the gap target.}
\label{tab:five-feedback-counts}
\begin{tabular}{llrrrrrr}
\toprule
Mode & Route & Queries & Accurate & Used & Late & Busy & Warnings\\
\midrule
Ideal & SDP & 48 & 42 & 48 & 41 & 0 & 39\\
Ideal & Dual & 48 & 48 & 48 & 7 & 0 & 0\\
One period & SDP & 17 & 2 & 0 & 17 & 31 & 17\\
One period & Dual & 45 & 45 & 42 & 3 & 3 & 0\\
\bottomrule
\end{tabular}
\end{table}

Most importantly, every SDP request in deadline mode is late, so no new
plan is activated. Five of those 17 candidates also exceed the command cap;
they are not clipped and relabeled feasible. The command stays zero throughout
all four deadline-SDP episodes. Nominal positions therefore remain unchanged,
and disturbances produce substantial drift. The dual activates 42 plans,
retaining three late discards and three occupied-worker skips. These outcomes
do not support a real-time advantage for the present SDP implementation;
they do not show that every possible SDP implementation must fail.

Table~\ref{tab:five-feedback-errors} reports every scenario and mode.
Position RMS is computed from the exact continuous integral of squared
position over all axes and six seconds, not only sampled positions. Full
terminal errors in all five derivatives and sampled RMS are also archived.
\begin{table}[htbp]
\centering\small
\caption{All fifth-order feedback outcomes, rounded for display. ``Delay''
denotes the fixed one-period measured-service-time replay. Small final position
does not imply that all higher derivatives have converged to zero.}
\label{tab:five-feedback-errors}
\begin{tabular}{llrrrr}
\toprule
Scenario & Mode & RMS, SDP & RMS, dual &
\shortstack{Final $|p|_\infty$\\SDP} & \shortstack{Final $|p|_\infty$\\Dual}\\
\midrule
One axis nominal & Ideal & 0.0456432 & 0.0456432 & 0.00130867 & 0.00130830\\
One axis nominal & Delay & 0.1000000 & 0.0539991 & 0.10000000 & 0.00327958\\
Two axes nominal & Ideal & 0.0417026 & 0.0411113 & 0.00564300 & 0.00130828\\
Two axes nominal & Delay & 0.0905539 & 0.0487769 & 0.10000000 & 0.00344047\\
Two axes pulse & Ideal & 0.0417272 & 0.0411504 & 0.00487318 & 0.00487282\\
Two axes pulse & Delay & 0.2215350 & 0.0496875 & 0.76761719 & 0.00072119\\
Two axes alternating & Ideal & 0.0416444 & 0.0411858 & 0.00407946 & 0.00340274\\
Two axes alternating & Delay & 0.2567869 & 0.0498793 & 0.97500000 & 0.02171400\\
\bottomrule
\end{tabular}
\end{table}

Ideal single-axis trajectories nearly coincide; two-axis paths can differ
because minimum-peak inputs need not be unique. The largest service time
is 1.752 seconds (rounded upward). Ideal execution does not make those
plans available in real time.

\subsubsection{Fourth-order instantaneous-planning comparison}\label{sec:four-feedback}
The earlier fourth-order test has four analogous scenarios, two methods and
20 replans per episode: eight episodes and 160 queries, with period $1/4$,
horizon two, command cap five and quantization $q=10^{-9}$.
Initial positions are $0.1,-0.08$, with other derivatives zero.
The pulse is $(0.05,-0.04)$ at steps 4--7 (zero-based); alternating
amplitudes are $1/30,1/40$, with four-step signs and a two-step axis offset.
Each planner sees its current state but not future disturbances.
Every eligible plan executes its full quantized prefix; absent eligibility,
the specified zero input is not claimed safe.

Unlike the fifth-order deadline mode, this simulation assumes instantaneous
planning. An independent rational convolution checker verifies states and
the local estimate
\begin{equation}\label{eq:feedback-local}
 |x^{\rm actual}_{a,d}(\delta)-x^{\rm nominal}_{a,d}(\delta)|
 \le (|w_a|+q/2)\frac{\delta^{4-d}}{(4-d)!},
\end{equation}
for constant step disturbance $w_a$ and $d=0,1,2,3$.
All 160 plans are eligible and meet \eqref{eq:tolerance}; maximum normalized
gaps are $9.924\cdot10^{-7}$ and $5.987\cdot10^{-7}$.
There are 337 optimizer calls, including 177 recovery LPs; SDP has 28
inaccurate statuses, whereas all 80 dual calls report success.
Table~\ref{tab:feedback} retains every scenario.
\begin{table}[ht]
\centering\small
\caption{Offline feedback results. Position RMS averages the initial and 20
step-end samples over all axes; final $|p|$ is the largest absolute axis
position at five seconds. Timing includes proposal, recovery and plan checking.
``Late'' means the measured planning time exceeds the simulated 0.25-second
period; the plant simulation does not incorporate this delay.}\label{tab:feedback}
\begin{tabular}{@{}llrrrr@{}}
\toprule
Scenario & Method & Pos. RMS & Final $|p|$ & Median s & Late/20\\
\midrule
Single nominal & SDP & 0.04319447 & 0.00264005 & 0.1279 & 0\\
Single nominal & Dual & 0.04319448 & 0.00264004 & 0.1052 & 0\\
Pair nominal & SDP & 0.03897921 & 0.00264002 & 0.2236 & 5\\
Pair nominal & Dual & 0.03893507 & 0.00264004 & 0.1991 & 0\\
Pair pulse & SDP & 0.03897522 & 0.00342167 & 0.2267 & 7\\
Pair pulse & Dual & 0.03895618 & 0.00342169 & 0.1985 & 0\\
Pair persistent & SDP & 0.03897294 & 0.00142946 & 0.2251 & 7\\
Pair persistent & Dual & 0.03895790 & 0.00146956 & 0.1976 & 0\\
\bottomrule
\end{tabular}
\end{table}

Command peaks range from 2.39999750 to 2.40000187, and no fallback occurs.
Nineteen of 80 SDP requests exceed the period, with maximum 0.667594 seconds;
the dual has none in this run, with maximum 0.224051 seconds.
Ignoring these delays prevents a real-time conclusion. Neither the local
error estimate nor nominal terminal equalities imply stability, robust
recursive feasibility or convergence of the disturbed plant.


\subsection{Conclusions and limitations}\label{app:reproducibility}
The explicit complete moment-body lift, its zero-peak homogenization and
segment concatenation provide the finite minimum-peak formulation. The
extension-degree bounds locate its representation class without asserting
minimum block count or minimum auxiliary dimension.
The exact outer model and its finite recovery algorithm must be distinguished:
a certified wide bracket remains feasible but is not converged, while a late
feasible plan remains unavailable at its deadline. The data support an
auditable offline formulation, not a general speed, stability or safety claim.
Both routes share certification, so their runtime ratio is not an isolated
price of certification. The results do not settle the exact extension degree
of $\K_5$, SOC representability of $\K_4$, or finite SDP representability
in higher orders. Reducing model size or improving deadline performance
requires further mathematical or computational evidence.

The reproducibility archive (\path{ESM_1.zip}) supplies the experimental protocols,
source code and recorded results, including warnings and failures. Its
reproducibility guide maps the claims to executable checks of the displayed
polynomial identities, rational certificates and feedback records. These
checks supplement, rather than replace, the analytic proofs.

\paragraph{Reproducibility materials.}
Archive contents: Source code, numerical records,
rational verification scripts and a reproducibility guide for the exact
semidefinite formulation and its computational studies, including warnings,
failed cases and boundary diagnostics.
The archive (\path{ESM_1.zip}, version 2026-10-09) is available from the
\href{https://github.com/moyuye77-code/Minimum-Peak-Control-of-Fifth-Order-Integrators-via-Exact-Semidefinite-Lifts/releases/tag/reproducibility-2026-10-09}{versioned GitHub release}.\footnote{Repository:
\begingroup\def\UrlBreaks{\do\/\do-\do.}
\url{https://github.com/moyuye77-code/Minimum-Peak-Control-of-Fifth-Order-Integrators-via-Exact-Semidefinite-Lifts}\endgroup.
Release tag: \texttt{reproducibility-2026-10-09}.}

\backmatter

\begin{thebibliography}{11}
\bibitem{averkov} Averkov, G.: Optimal Size of Linear Matrix Inequalities
in Semidefinite Approaches to Polynomial Optimization.
\emph{SIAM Journal on Applied Algebra and Geometry} \textbf{3}(1),
128--151 (2019).
\url{https://doi.org/10.1137/18M1201342}.
Author version: \url{https://arxiv.org/abs/1806.08656v2}.
\bibitem{fawzi} Fawzi, H.: On representing the positive semidefinite cone
using the second-order cone. \emph{Mathematical Programming}
\textbf{175}(1--2), 109--118 (2019).
\url{https://doi.org/10.1007/s10107-018-1233-0}.
\bibitem{gr} Gosse, L., Runborg, O.: Resolution of the finite Markov
moment problem. \emph{Comptes Rendus Mathematique} \textbf{341}(12),
775--780 (2005).
\url{https://doi.org/10.1016/j.crma.2005.10.009}.
Author version: \url{https://arxiv.org/abs/0910.5791v1}.
\bibitem{gn} Gouveia, J., Netzer, T.: Positive Polynomials and
Projections of Spectrahedra. \emph{SIAM Journal on Optimization}
\textbf{21}(3), 960--976 (2011).
\url{https://doi.org/10.1137/100801913}.
Author version: \url{https://arxiv.org/abs/0911.2750v2}.
\bibitem{gpt} Gouveia, J., Parrilo, P.A., Thomas, R.R.: Lifts of Convex
Sets and Cone Factorizations. \emph{Mathematics of Operations Research}
\textbf{38}(2), 248--264 (2013).
\url{https://doi.org/10.1287/moor.1120.0575}.
Author version: \url{https://arxiv.org/abs/1111.3164v2}.
\bibitem{intersection} Haddad, S., Halder, A.: Certifying the Intersection
of Reach Sets of Integrator Agents with Set-valued Input Uncertainties.
\emph{IEEE Control Systems Letters} \textbf{6}, 2852--2857 (2022).
\url{https://doi.org/10.1109/LCSYS.2022.3179666}.
Author version: \url{https://arxiv.org/abs/2203.12007v4}.
\bibitem{hh} Haddad, S., Halder, A.: The Curious Case of Integrator Reach
Sets, Part I: Basic Theory. \emph{IEEE Transactions on Automatic Control}
\textbf{68}(11), 6680--6695 (2023).
\url{https://doi.org/10.1109/TAC.2023.3244694}.
Author version: \url{https://arxiv.org/abs/2102.11423v2}.
\bibitem{lasserredensity} Lasserre, J.B.: Borel measures with a density on a compact
semi-algebraic set. \emph{Archiv der Mathematik} \textbf{101}(4),
361--371 (2013).
\url{https://doi.org/10.1007/s00013-013-0557-5}.
Author version: \url{https://arxiv.org/abs/1304.1716v2}.
\bibitem{mcclure} McClure, D.E.: Perfect spline solutions of $L_\infty$
extremal problems by control methods. \emph{Journal of Approximation Theory}
\textbf{15}(3), 226--242 (1975).
\url{https://doi.org/10.1016/0021-9045(75)90105-7}.
\bibitem{nie} Nie, J.: Polynomial Matrix Inequality and Semidefinite
Representation. \emph{Mathematics of Operations Research}
\textbf{36}(3), 398--415 (2011).
\url{https://doi.org/10.1287/moor.1110.0498}.
Author version: \url{https://arxiv.org/abs/0908.0364v2}.
\bibitem{nierational} Nie, J.: First Order Conditions for Semidefinite
Representations of Convex Sets Defined by Rational or Singular Polynomials.
\emph{Mathematical Programming} \textbf{131}, 1--36 (2012).
\url{https://doi.org/10.1007/s10107-009-0339-9}.
Author versions: \url{https://arxiv.org/abs/0806.4721v1} and the
July 2009 revision at
\url{https://mathweb.ucsd.edu/~njw/PUBLICPAPERS/QmPoSdpR_mpa2.pdf}.
\end{thebibliography}

\section*{Statements and Declarations}
\bmhead{Funding}
This research was self-funded by the author. No external funding was received.

\bmhead{Competing interests}
The author declares no financial or non-financial competing interests
relevant to this work.

\bmhead{Use of artificial intelligence}
The author determined the research direction and prepared preliminary
derivations and text. OpenAI Codex assisted with checking boundary cases,
refining mathematical arguments, developing simulation and verification
code, analyzing results, and drafting and editing the manuscript. The
resulting derivations, computational results and manuscript were
subsequently reviewed by the author.

\bmhead{Data and code availability}
The source code, frozen numerical records and rational verification scripts
supporting this study are publicly available in the
\href{https://github.com/moyuye77-code/Minimum-Peak-Control-of-Fifth-Order-Integrators-via-Exact-Semidefinite-Lifts/releases/tag/reproducibility-2026-10-09}{GitHub reproducibility archive},
version 2026-10-09 (\path{ESM_1.zip}).
\end{document}
